\documentclass[ekobook.tex]{subfiles}

\begin{document}\setWPP
\lsection{Obchodní korporace}
\qaw{Co je to obchodní korporace, kdo je FO a PO}{%
Obchodní korporace je \texthl{právnická osoba} založená jednou nebo více osobami za účelem podnikání nebo jiného cíle.\\
\textbl{FO (fyzická osoba)} -- konkrétní člověk s právy a povinnostmi\\
\textbl{PO (právnická osoba)} -- organizovaný útvar (společenství osob), který má vlastní právní osobnost
}
\qaw{Kdo jedná jménem společnosti (FO, PO, prokura)}{%
\textbl{Fyzická osoba}
-- jedná \texthl{osobně} nebo prostřednictvím \texthl{zplnomocněného zástupce}\\
\textbl{Právnická osoba}
-- jedná prostřednictvím \texthl{statutárního orgánu} (jednatel, představenstvo, správní rada) nebo zástupce\\
\textbl{Prokura}
-- zvláštní forma plné moci; prokurista je oprávněn ke \texthl{všem právním úkonům} týkajícím se podniku; musí být zapsán v obchodním rejstříku
}
\qaw{Čím se řídí obchodní korporace}{%
Obchodní korporace se řídí především \texthl{Zákonem o obchodních korporacích} (zák. č. 90/2012 Sb.) a \texthl{Občanským zákoníkem}.
Dále stanovami (nebo společenskou smlouvou) a vnitřními předpisy společnosti.
}
\qaw{Jak se zakládá OK}{%
Obchodní korporace se zakládá \texthl{společenskou smlouvou} (při více zakladatelích) nebo \texthl{zakladatelskou listinou} (při jednom zakladateli). Oba dokumenty musí mít formu \texthl{notářského zápisu}.\\
Společnost \texthl{vzniká} dnem \texthl{zápisu do obchodního rejstříku} (vede ho rejstříkový soud). Obchodní rejstřík je veřejný seznam.
}
\qaw{Vysvětli základní rozdíl mezi osobními a kapitálovými společnostmi -- ručení}{%
\textbl{Osobní společnosti} (v.o.s., k.s.)
-- společníci ručí \texthl{neomezeně} (celým svým majetkem), často solidárně; osobní účast na řízení; nemají povinný minimální základní kapitál\\
\textbl{Kapitálové společnosti} (s.r.o., a.s.)
-- ručení je \texthl{omezené} nebo žádné (společníci/akcionáři neručí osobním majetkem); oddělení majetku společnosti od majetku společníků; povinný základní kapitál
}
\qaw{Popiš k.s., v.o.s.:}{%
\textbl{v.o.s. (veřejná obchodní společnost)}
-- minimálně 2 společníci; ručí \texthl{neomezeně a solidárně} celým svým majetkem; zisk a ztráta se dělí rovným dílem (nebo dle smlouvy);
každý společník je statutárním orgánem\\
\textbl{k.s. (komanditní společnost)}
-- dva typy společníků: \texthl{komplementář} (ručí neomezeně, řídí společnost) a \texthl{komanditista} (ručí omezeně do výše vkladu, neřídí);
zisk se dělí obvykle na polovinu nebo dle smlouvy
}
\qaw{Popiš s.r.o., a.s.:}{%
\textbl{s.r.o. (společnost s ručením omezeným)}
-- minimální vklad 1 Kč; společníci ručí \texthl{omezeně} do výše nesplaceného vkladu;
může založit i jedna osoba; statutární orgán = jednatel(é); nejvyšší orgán = valná hromada\\
\textbl{a.s. (akciová společnost)}
-- minimální základní kapitál 2 000 000 Kč (nebo 80 000 EUR);
akcionáři \texthl{neručí} za závazky společnosti; základní kapitál rozdělen na akcie
}
\qaw{a.s.:}{%
\textbl{monistický systém}
-- statutární orgán = \texthl{správní rada} + statutární ředitel\\
\textbl{dualistický systém}
-- statutární orgán = \texthl{představenstvo}, kontrolní orgán = \texthl{dozorčí rada}\\
\textbl{dividenda}
-- podíl na zisku vyplácený akcionářům\\
\textbl{zatímní list}
-- doklad vydávaný před splacením akcií (nahrazuje akcie do doby splacení)\\
\textbl{tantiéma}
-- odměna členům představenstva/dozorčí rady/spravní rady z zisku\\
\textbl{valná hromada}
-- nejvyšší orgán společnosti (rozhoduje o zásadních otázkách)
}
\qaw{Rozdíl mezi listinnou a zaknihovanou akcií}{%
\textbl{Listinná akcie}
-- existuje jako \texthl{fyzický cenný papír} vytištěný na papíře s ochrannými prvky\\
\textbl{Zaknihovaná akcie}
-- existuje pouze jako \texthl{elektronický záznam} v Centrálním depozitáři cenných papírů (od 2014 jsou akcie na majitele povinně zaknihované)
}
\qaw{Vyjmenuj druhy akcií}{%
\textbl{kmenové}
-- standardní akcie s hlasovacím právem a právem na dividendu\\
\textbl{prioritní}
-- přednostní výplata dividendy, obvykle omezené nebo žádné hlasovací právo\\
\textbl{zaměstnanecké}
-- vydávány zaměstnancům jako motivace\\
\textbl{na jméno} / \textbl{na majitele}
-- podle způsobu převodu
}
\wpage
\end{document}